
\begin{abstract}

蔬菜类商品保鲜期短、品相随时间变差，商超需每日做出补货与定价决策。本文基于某商超2020年7月至2023年6月六个蔬菜品类251个单品的销售流水、批发价格和损耗率数据，建立了从品类关联分析、品类级补货定价到单品级选品优化的递进式数学模型。

\textbf{对于问题一,} 首先对销售流水按品类和日期聚合，经正态性检验发现全部品类拒绝正态分布假设，采用Spearman秩相关系数分析品类间关联关系，接着使用K-means++聚类算法对159个销售稳定的单品按销量规模和波动性聚类，最后对各品类日销售量进行分布拟合。结果表明，Spearman $\rho$ 范围为$[-0.194, 0.625]$，花叶类与花菜类关联最强（$\rho=0.625$），单品被聚为低量高波动与高量中波动两类（Silhouette=0.463），对数正态分布最能刻画蔬菜日销售量的分布形态。

\textbf{对于问题二,} 首先建立Prophet时间序列分解模型分离趋势、周/年季节性和中国节假日效应，接着从数据自行估计各品类需求价格弹性，然后建立约束报童模型以加成率和订货量为联合决策变量，通过二维网格搜索求解最优补货与定价策略。结果表明，单日总期望利润为2,918元（折扣回收率$k=0.66$，$k\in[0.50,0.75]$对应区间$[2,362,\;3,713]$元），较ARIMA+固定加成baseline（723元/天）提升304\%。优化加成率在103.6\%--128.2\%之间，品类间有充分区分。

\textbf{对于问题三,} 首先从6月24--30日可售品种中过滤日均销量$\geq 2.5$ kg的单品得到31个（恰好落在27--33范围内），接着各单品独立运行报童优化确定订货量与定价。结果表明，单日总期望利润为1,327元，较固定加成baseline（465元）提升185\%。全部六个品类需求满足率均超过78\%，水生根茎类最低（78.6\%），其余品类超过150\%。

\textbf{对于问题四,} 基于前三问建模过程中实际暴露的数据局限性，建议商超优先采集每日库存与缺货记录、折扣原因分类码和天气数据，以提高最敏感参数$k$的估计精度和需求预测质量。模型具有清晰的数学结构和可解释的参数，关键数值结论均经过参数扰动和敏感性检验。

\keywords{蔬菜定价；补货决策；Spearman秩相关；Prophet时间序列分解；报童模型；K-means++聚类}

\end{abstract}
